\documentclass[12pt]{scrartcl}
\usepackage[sexy, href, sectionmark, mdthm]{moksh}
\setlength{\parindent}{0pt}
\setlength{\parskip}{8pt}
\DeclareMathOperator{\ob}{ob}


\makeatletter
\let\c@exercise\c@theorem
\makeatother
% makes exercises the same counter as theorem/everything else.

\newcounter{exercises}[section]
\renewcommand\theexercises{\thesection.\arabic{exercises}}
\newenvironment{exercises}[1][]
{\refstepcounter{exercises}\par\medskip\noindent \textbf{Exercise \theexercises}\ifx#1\empty\else\space(#1)\fi\par} {\par\medskip}

	\newcounter{soln}[section]
\renewcommand{\thesoln}{\thesection.\arabic{soln}}
\renewenvironment{soln}
  {\refstepcounter{soln}%
   \begin{mdframed}[style=mdplainbox]\noindent\textbf{Solution \thesoln:}\par\smallskip\ignorespacesafterend}
  {\end{mdframed}}



\usepackage[all,2cell]{xy}
\UseAllTwocells
\xyoption{arc}
\xyoption{rotate}
\usepackage{enumitem}


\title{225A: Topics in Number Theory}
\author{moksh}

\begin{document}

\maketitle
\tableofcontents

\newpage

Topics we plan to cover

\begin{enumerate}
	\item Galois Cohomology
	\item Central simple algebras
	\item Brauer Groups and its relation with Galois Cohomology
	\item Brauer group of local and global fields
\end{enumerate}


\begin{motivation}[Group Cohomology]
Let $G$ be a discrete topological group. (topology is discrete)
There exists a \vocab{classifying space} of $G$, written as $BG$.
Is an Eilenberg-Maclane space, $K(G, 1)$.
In other words, 
$	\pi_i(BG) = 
\begin{cases}
	G & \textrm{if } i=1  \\
	0 & \textrm{if } i>1  \\

\end{cases}
$.

Existence can be shown by categorical means, has the universal property that... 

Classifies all possible $G$-bundles.

What about singular cohomology? $H_{\textrm{sing}}^{*} (BG, \mathbb{Z})$ 

Usually homotopy groups are more complicated than homology groups, but the reverse is the case here.
Very tough in general to compute the homology groups, and all homotopy groups are known.
\end{motivation}
This computation of singular homology groups is a good motivation for Group Cohomology, gives us a lot of information about the group $G$ itself.
	
Whole idea of Rep. Theory is to study the group itself, same idea here. 
Calculate lots of invariants of $G$, such as information about the classifying space.
For example, if $G$ is a Galois group, then its group cohomology gives us a lot of information.

\begin{example}
Let $G = \mathbb{Z}$.
We have that $BG \cong S^1$.
We have the universal property
\[
\begin{array}{c}
	\xymatrix{
		\mathbb{R}\ar[d] \\
S^1
	}
\end{array}
\]
$H^{*}(S^1, \mathbb{Z})$ is well known.
\end{example}

\begin{example}
$G = \mathbb{Z}/2 \mathbb{Z}$.
$BG \cong \mathbb{RP}^{\infty}$.
\[
\begin{array}{c}
	\xymatrix{
	S^{\infty}\ar[d] \\
\mathbb{RP}^{\infty}
	}
\end{array}
\]
Hence, $H^{*}(S^1, \mathbb{Z}) \cong H^{*}(\mathbb{RP}^{\infty}, \mathbb{Z})$.
\end{example}
These are what are called \vocab{Lens spaces}, and as we can see, can get complicated fast, since its already $\mathbb{RP}^{\infty}$ for $\mathbb{Z}/2 \mathbb{Z}$, we haven't even gotten to $\mathbb{Z}/p \mathbb{Z}$.
Hence, studying this through a topological lens is quite difficult, which is why do things in a very algebraic manner.

Group cohomology is isomorphic to the singular cohomology of $BG$, and its far easier to compute. 

Let $G$ be a finite group.
This isn't a necessary hypothesis, we just need that $G$ has the discrete topology. 
However, we often work with Galois groups of finite extensions.
This also helps us in the infinite case, because infinite Galois groups arise as inverse limits of all finite groups, with the profinite topology.
\begin{definition}
Let $A$ be an abelian group with a $G$-action.
In other words, there exists a homomorphism $\alpha: G \to \Aut(A)$.
In this case, we say that $A$ is a $G$-module.
Let $\catname{G-mod}$ be the category of all $G$-modules
\end{definition}

\begin{exercise}
	$\catname{G-mod}$ is an abelian category. 
\end{exercise}

Note that $\alpha$ can be the trivial homomorphism, so any abelian group can be represented as a $G$-module, so $\catname{G-mod}$ is non-trivial.
Let $\mathbb{Z}[G]$ be the group algebra of $G$ over $\mathbb{Z}$.
\begin{definition}
	$\mathbb{Z}[G]$, as an abelian group, is the free abelian group on $G$, can be thought of as all $\mathbb{Z}$-linear combinations of $G$, with multiplication in the obvious manner.
\end{definition}
	\begin{itemize}
		\item $\mathbb{Z}[G]$ is abelian $\iff$ $G$ is abelian.
		\item $\mathbb{Z}[G]$ is left-Noetherian $\iff$ $G$ is finitely generated.
	\end{itemize}
	There exists a $\mathbb{Z}$-algebra map
	\begin{align*}
		\epsilon: \mathbb{Z}[G] &\to \mathbb{Z} \\
		\epsilon \left(\sum m_i g_i\right) &= \sum m_i
	\end{align*}
	such that
	\[
		\begin{array}{c}
			\xymatrix{
				\mathbb{Z} \ar[r]\ar[dr]_{\textrm{id}} & \mathbb{Z}[G] \ar[d]^{\epsilon} \\
							 & \mathbb{Z} \\
						 }	
	\end{array}
	\]
$\mathbb{Z}[G]$ is an \vocab{augmented} $\mathbb{Z}$-algebra.
Let $I_G = \ker(\epsilon)$, which is the \vocab{augmented ideal of} $\mathbb{Z}[G]$.
It is a two-sided ideal, and $\mathbb{Z} = \mathbb{Z} \oplus I_G$.

\begin{lemma}
	$I_G$ is the free abelian group on the set \textcolor{red}{finish later}
\end{lemma}

Let $G$ be a group, and let $A$ be a $G$-module.
Then, $A$ is canonically a $\mathbb{Z}[G]$-module with
\[
	\left(\sum m_i g_i\right) a = \sum m_i  \left(g_i a\right).
\]
Conversely, any module over $\mathbb{Z}[G]$ is a $G$-module.

\begin{proof}
There exists a group homomorphism $\Phi: G \injto \mathbb{Z}[G]^X$ by sending $g$ to $1 \cdot g$.
Similarly, there exists a map $\mathbb{Z}[G]^{\times }\to \Aut(A)$.
The composition is a group homormophism from $G$ to $\Aut(A)$, which is the definition of a group action, so we are done.
\end{proof}

Remember that for a fixed ring $R$, we have that category $\catname{R-mod}$ (left $R$-modules) is an abelian category. 
\begin{proposition}
	There exists an equivalence of abelian categories $\catname{G-mod} \to \catname{\mathbb{Z}[G]-mod}$ 
\end{proposition}
\begin{proof}
The previously discussed correspondence.
\end{proof}

Thus, we make no distinction between these two categories in this class.

There is a functor 
\begin{align*}
	F: \catname{G-mod} &\to \catname{Ab} \\
	A  &\mapsto A^G
\end{align*}
where $A^G$ is the subgroup of $G$-invariant elements of $A$.
We now wonder properties this functor has, left-exact, right-exact, exact?

\begin{fact}
Consider the trivial functor that represents an arbitrary abelian group as a trivial $G$-module.
Then, $\left( \left(\dashbk\right)^{tr}, \left(\dashbk\right)^{G}\right)$ are an adjoint pair, and we have that
$\Hom_{\catname{Ab}}\left(A, B^G\right) \to_{\sim} \Hom_{\catname{G-mod}} \left(A^{tr}, B\right)$ 
\end{fact}

\begin{corollary}
 $\left(\dashbk\right)^G$ is left exact. 
\end{corollary}

We have that for an exact sequence
\[
0 \to A' \to A \to A'' \to 0
\] 
in $\catname{G}-mod$,
\[
0 \to A'^G \to A^G \to A''^G 
\] 
is exact.

\begin{exercise}
Showcase an instance where the last map is not surjective.
\end{exercise}
Hence, this functor is left-exact but not necessarily right exact.
In homological algebra, when we have such a situation, we construct the right-derived functor, and this is what we call Group Cohomology.
The existence of this right-derived functor happens when we have enough injectives, i.e. every object injects inside an injective object.

\section{Class 2}
\begin{theorem}
For an arbitrary ring $R$, the category $\catname{R-mod}$ has enough injectives.
\end{theorem}
\begin{proof}
	(Sketch)
	\begin{lemma}
	Let $A$ be an abelian group.
	Then, there exists an injective abelian group $I$ and an injection of abelian groups $A \injto I$  
	\end{lemma}
	\begin{proof}
	Let $S = \Hom_{\catname{Ab}}(A, \mathbb{Q}/\mathbb{Z})$, and $I = \left(\mathbb{Q}/\mathbb{Z}\right)^S$.
	Define $A \to I$, by sending every element $a$ to $ \left(\alpha(a)\right)_{\alpha \in S}$.
	Check that this has no kernel, since there is no element that sends to $0$ under every homomorphism into $\left(\mathbb{Q}/\mathbb{Z}\right)$  
	\end{proof}
\[
	\begin{array}{c}
\xymatrix{
	M \ar[r]_{\cong} \ar[d]& \Hom_R(R, M) \ar@{^{(}->}[d] \\
	\Hom_{\mathbb{Z}}(R, I) & \Hom_{\mathbb{Z}}(R, M) \ar@{^{(}->}[l] \\
}
\end{array}
\] 
where $M \injto I$ as in the \textcolor{red}{the lemma}. 
Thus, it is enough to show that $I' = \Hom_{\mathbb{Z}}(R, I)$ is an injective $R$-module.

\begin{proposition}
Let $A \injto B$ be an inclusion of $R$-modules.
We need to show that $\Hom_R(B, I') \surjto \Hom_R(A, I')$ by extension of scalars.

\[
	\begin{array}{c}
\xymatrix{
	\Hom_R(B, I') \ar[r] \ar[d]_{S} & \Hom_R(A, I') \ar[d]^{S} \\
	\Hom_{\mathbb{Z}}(B, I) \ar@{->>}[r]& \Hom_{\mathbb{Z}}(A, I) \\
}
\end{array}
\] 
\end{proposition}
\end{proof}
\begin{corollary}
	$\catname{\mathbb{Z}[G]-mod}$ has enough injectives 
\end{corollary}

\begin{lemma}
	$A^G = \Hom_{\mathbb{Z}[G]}(\mathbb{Z}, A)$
\end{lemma}
\begin{proof}
	By adjointness, $\Hom_{\mathbb{Z}[G]}(\mathbb{Z}, A) \to_{\cong} \Hom_{\mathbb{Z}}(\mathbb{Z}, A^G) = A^G$.
\end{proof}
Let $A \to I$ be an injective resolution in $\catname{\mathbb{Z}[G]-mod}$.
Then, we define $\Ext_{\mathbb{Z}[G]}^{i} (\mathbb{Z}, A) = H^i \left(\Hom_{\mathbb{Z}[G]}(\mathbb{Z}, I)\right)$, for $i \geq 0$.

\begin{definition}
	We let $H^{i}(G, A) = \Ext_{\mathbb{Z}[G]}(\mathbb{Z}, A)$ and call this the $i$'th group cohomology of $G$ with coefficients in $A$.
\end{definition}

\begin{corollary}
$H^0(G, A) \cong A^G$ 
\end{corollary}

\begin{theorem}
$\Ext_R^i(M, N) = 0$ for all $i>0$ if $M$ is projective (also true for $M$ injective).   
\end{theorem}

Thus, we can take a projective resolution of $M$, or injective resolution of $N$ to compute the group cohomology.
Then, if we can find a free resolution of $\mathbb{Z}$ over $\mathbb{Z}[G]$, we are done.
\textbf{This is the central idea of group/Galois cohomology, to find some canonical free/projective resolutions of $\mathbb{Z}$.}

\begin{corollary}
	If $f: A \to B$ is a $G$-equivariant map (morphism in $G\catname{-mod}$), then there exists a canonical map $f^i: h^i(G, A) \to H^i(G, B)$. 
\end{corollary}

\begin{corollary}
Recall that $H^0(G, A) = A^G$.
We also have that $H^1(G, A)$ is the group of equivalence classes of extensions of $\mathbb{Z}$ by $A$ in $\catname{\mathbb{Z}[G]-mod}$  
\end{corollary}
An \vocab{extension} of $A$ by $B$ is an exact sequence
	\[
	0 \to B \to H \to A \to 0
	\] 
	Two exactensions are said to be isomorphic iff there is an isomorphism $H \to H'$ such that the morphism of exact sequences commutes.
	Equivalence classes of extensions admit a group action, given by the fiber product.
	\begin{corollary}
	Let $0 \to A' \to A \to A''$ be a an exact sequence of $G$-modules.
	Then, there exists a long exact sequence
	\[
\dots \to H^i(G, A') \to H^i(G, A) \to H^i(G, A'') \to H^{i+1}(G, A^i) \to \dots	
	\] 
	\end{corollary}
	
	Let $G$ be a finite group (not necessarily abelian).
	We define an element $N \in \mathbb{Z}[G]$ as $N = \sum_{g \in G} g$.
	This is called the \vocab{norm element} of $\mathbb{Z}[G]$, related to the norm of an element in the context of Galois theory. 

	\begin{proposition}
		$N \in Z(\mathbb{Z}[G])$. 
	\end{proposition}
	\begin{proof}
	Check it.
	\end{proof}
	\begin{proposition}
		$N$ is $G$-invariant under the action of $G$ on $\mathbb{Z}[G]$. 
	\end{proposition}
	\begin{proof}
	Trivial.
	\end{proof}

	\begin{theorem}
	Let $|G| = m < \infty$.
	Let $R$ be a commutative ring with $m \in R^x$, and let $A$ be a $G$-module s.t. $A$ is also an $R$-module.
	Then, $\frac{N}{m}$ is a central idempotent in $R[G]$.
	Moreover, $H^i(G, A) = 0$ for all $i \neq 0$.  
	\end{theorem}
	This shows that in this situation, the theory of Galois cohomology is trivial.
	For instance if $A$ is a $\mathbb{Q}$ or $\mathbb{R}$ or $\mathbb{C}$ vector space, it gives you no information.
	This is because Galois cohomology is a \emph{torsion phenomenon}.
	\begin{proof}
	It is easy to see that $\frac{N}{m} $ is central.
	Note that $N^2 = N \left(\sum g_i\right) = \sum g_i N$ since $N$ is central.
	However, this is equal to $mN$, so $\left(\frac{N}{m}\right)^2 = \frac{N}{m}$
	Let $e$ represent $\frac{N}{m}$.

	From Atiyah and Macdonald Ch.1 exercises, recall that a nontrivial idempotent splits a ring into the direct product of subrings; we can write $R[G] = R[G]e \times R[G]e'$ where $e' = 1-e$.
	Let these be represented by $R_1, R_2$ respectively.
	Since these are both direct summands of $R[G]$, which is free, we have that they are both projective modules.


	Let $ $ 
	Thus, we have that $\Ext_{R[G]}^i(R_1, M) = 0$ for all $R[G]$-modules $m$ and every $i\neq 0$, by taking a projective resolution.
		We also have that $H^0(G, A) = eA$ 
	\end{proof}
	
	\section{Class 3}
	\begin{corollary}
		If $G$ is finite and $A$ is a $G$-module s.t. $A$ is also $\mathbb{Z}[1/m]$ module, e.g. a $\mathbb{Q}$-vector space, then $H^i(G, A) = 0$ for all $i \neq 0$.     
	\end{corollary}
	
		\begin{corollary}
			If $G$ is finite, then $H^0(G, \mathbb{Z}) = \mathbb{Z}$, $H^1(G, \mathbb{Z}) = 0$, and $H^i(G, \mathbb{Z}) \cong H^{i-1}(G, \mathbb{Q}/\mathbb{Z})$ for all $i>0$.   
		\end{corollary}

		\begin{proposition}
		We have an exact sequence of trivial $G$-modules
		\[
		0 \to \mathbb{Z} \to \mathbb{Q}\to \mathbb{Q}/\mathbb{Z} \to 0
		\] 
		\end{proposition}

		\begin{corollary}
		We have a long exact sequence
		\[
			0 \to H^0(G, \mathbb{Z}) \to H^0(G, \mathbb{Q}) \to H^0(G, \mathbb{Q}/\mathbb{Z}) \xrightarrow{0} H^1(G, \mathbb{Z}) \to \dots
		\] 
		\end{corollary}

		Let $G$ be any discrete group.
		We construct a free resolution of $\mathbb{Z}$ as a $\mathbb{Z}[G]$-module.
		Let $B_0 = \mathbb{Z}[G]$, and for $n>0$ let $B_n$ be the free $\mathbb{Z}[G]$-module on $G^n$   

		We notice that $\Hom_{\mathbb{Z}[G]}(B_n, B_{n-1}) = \Hom_{\catname{Set}}(G^n, B_{n-1})$.

		\textcolor{red}{fill in details from images i took}
		Let $d^n = \sum_{i=0}^n (-1)^{n} d_i^n: G^n \to B_{n-1}$. 
		
something something simplicial objects in categories, Dold-Kan correspondence between simplicial abelian groups and chain complexes.

		\begin{theorem}
		This sequence is exact
		\end{theorem}
		\begin{proof}
		
		
To show a complex is acyclic, we can manually compute boundary maps, show all homologies are zero.
But instead, we can prove the strong statement, which is that it is nullhomotopic, which is often easier.

We want to construct maps $h_n: B_n \to B_{n+1}$ for all $n \geq -1$ such that $h_{n-1} \circ d^n - d^{n+1} \circ h_n = \id_{B_n}$.
Define $h_n: B_n \to B_{n+1}$ 
		\end{proof}

		\section{Class 4}
Let $G$ be a discrete group.
The \vocab{bar resolution} of $\mathbb{Z}$ as a $\mathbb{Z}[G]$ module is the above resolution.
If $A$ is a $G$-module, then write $C^n(G, A) = \Hom_{\mathbb{Z}[G]}(B_n, A) = \Hom_{\catname{Set} }(G^n, A)$.

$Z^n(G, A) = \ker(d_n^{*})$

$B^n(G, A) = \img(d_{n-1}^{*})$.


$H^n(G, A) = Z^n(G, A)/B^n(G, A) = \frac{n-\textrm{cocycles in }A }{n-\textrm{coboundaries in }A} $ 

\begin{theorem}
	Suppose $G$ is finite of order $m$, and let $A$ be a $G$-module.
	Then, $H^n(G, A)$ is annihilated by $m$, i.e., it is a $\mathbb{Z}/m$-module for $n>0$. 
\end{theorem}

\begin{proof}

:wLet $B \xrightarrow{\varepsilon} \mathbb{Z}$ be the bar resolution of $\mathbb{Z}$ as a $\mathbb{Z}[G]$ module.

Consider the map of complexes

\[
\xymatrix{
	\dots \ar[r]^{d_2} \ar[d]_{m} & B_1 \ar[r]^{d_1}\ar[d]_{m}& B_0 \ar[d] _{m-N} \\
	\dots \ar[r]^{d_2} & B_1 \ar[r]^{d_1} & B_0 \\
}
\] 

It remains to show that this map is null-homotopic in the degrees.

Let $\nu_n: B_n \to B_{n+1}$ be the map given by $\nu_n(g_1 | \dots | g_n) = \epsilon_n \sum_{g \in G} (g_1 | \dots | g_n | g$
where $\epsilon_n = (-1)^{n+1}$.
We need to show that $\nu_{n-1} \circ d_n + d_{n+1} \circ \nu_n = m$ is the degree. 

\end{proof}

\begin{corollary}
	If $A$ is $\mathbb{Z}[1/m]$ module with $G$-action, then $H^i(G, A) = 0$ for all $i>0$.  
\end{corollary}
\begin{proof}
$m$ is an automorphism of $A$, so it is an automorphism of $H^*(G, A)$ by applying cohomology functor.
But $m$ annihilates $H^i(G, A)$, so if an automorphism annihilates the group, the group is trivial. 
\end{proof}

\begin{lemma}
Suppose $|G| = m$, and $A$ is a finitely generated abelian group that is a $G$-module over $A$.
Then, all cohomology groups are also finitely generated abelian.
\end{lemma}

\begin{corollary}
Let $|G| = m$ and let $A$ be a $G$-module where is a finitely generated abelian group.
Then, $H^0(G, A)$ is a finite generated abelian group and $H^i(G, A)$ is finite for $i>0$.  
\end{corollary}

\begin{proof}
By the lemma, $H^i(G, A)$ is finitely generated abelian.
But it is also killed by $m$, hence torsion, and finite.
\end{proof}
Let $G= \cyc{n}$.
Let $N \in \mathbb{Z}[G]$ be the norm element.
Note that $\mathbb{Z}[\mathbb{Z}] = \mathbb{Z}[t, t^{-1}]$, so $\mathbb{Z}[\cyc{n}] = \mathbb{Z}[t]/t^n-1$.
Let $G =\langle \sigma\rangle$.
Then, $\sigma = t$ and $N = 1 + \sigma + \sigma^2 + \dots + \sigma^{n-1}$.

\section{Class 6}

Let $L/K$ be a finite Galois extension with group $G$, then $H^1(G, L^{\times }) = 0$, as we saw last class.
What is $H^1(G, L)$?

\begin{lemma}
Suppose $G$ is a discrete group and $H \injto G$ is a subgroup of finite index. 
Then, under the inclusion of $\mathbb{Z}$-algebras $\mathbb{Z}[H] \injto \mathbb{Z}[G]$, we have that $\mathbb{Z}[G]$ is a free $\mathbb{Z}[H]$-module with basis $G/H$.
\end{lemma}
\begin{proof}
We knew that $\mathbb{Z}[H]$ is a free $\mathbb{Z}$-module with basis $H$ and $\mathbb{Z}[G]$ is a free $\mathbb{Z}$-module with basis $G$.
Let $G/H = \left\{He, Hg_1, \dots, Hg_n\right\}$.

Then, note that the map $\Phi: \bigoplus_{i=0}^n \mathbb{Z}[H] \ol{g_i} \to \mathbb{Z}[G]$ is an isomorphism.
\end{proof}
Let $\alpha: R \to S$ be a (not necessarily commutative) ring homomorphism.
Assume that $S$ is projective (e.g. free) $R$-module via $\alpha$.
	Let $M \in R-\catname{mod}$.
	Then, we can define two $S$-modules 
	\begin{enumerate}[label=\arabic*)]
		\item $M\otimes_R S$.
		\item $\Hom_R(S, M)$. 
	\end{enumerate}
\begin{proposition}
$M \otimes_R(S)$ is called the induced $S$-module.
$\Hom_R(S, M)$ is called the coinduced $S$-module. 
Under these conditions, a canonical isomorphism $\Ext_S^i(N, \Hom_R(S, M)) \cong \Ext_R^i(N, M)$ exists for all $N \in S-\catname{mod}$ and $i\geq 0$.   
\end{proposition}

\begin{proposition}
Let $\alpha: R \to S$ be a (not necessarily commutative) ring homomorphism.
Assume that $S$ is projective (e.g. free) $R$-module via $\alpha$.
Let $P, M \in S-\catname{mod}$, $N \in R-\catname{mod}$, and $h: M \to N$ an $R-\catname{mod}$ homomorphism.

Under these conditions, there is a canonical isomorphism $\Ext_S^i(P, \Hom_R(S, M)) \cong \Ext_R^i(P, N)$. 
\end{proposition}

\begin{proof}
Let $P^{\bullet} \to N$  be a free $S$-module resolution.
Then, $P^{\bullet} \to N$ is also a projective (free, if $S$ is free over $R$) $R$-module resolution.
\begin{align*}
\Rightarrow \Ext_R^i(N, M) &= H^i(\Hom^{\bullet}_R(P^{\bullet}, M)) \\
 &= H^i(\Hom_S(P, \Hom_R(S, M))) \\
&= \Ext_S^i(N, \Hom_R(S, M)) \\
\end{align*}
\end{proof}

\begin{corollary}[Shapiro Lemma]
Let $G$ be a discrete group and $H \injto G$  a subgroup of finite index.

Take $R = \mathbb{Z}[H]$, $S= \mathbb{Z}[G]$.
Let our $M$ in the previous prop. refer to any $H$-module, $A$ and let our $N$ refer to $\mathbb{Z}$ as a trivial $G$-module.
\end{corollary}
This satisfies the hypotheses, so we get that $\Ext_S^i(\mathbb{Z}, \textrm{coind}_H^G(A)) \cong \Ext_R^i(\mathbb{Z}, A) \Leftrightarrow H^i(G, \textrm{coind}_H^G(A)) \cong H^i(H, A)$ for all $i$. 

\begin{corollary}
Let $G$ be a finite group and let $A$ be any abelian group.
Then, $H^i(G, \textrm{coind}_{*}^G(A)) = \begin{cases} A & \text{if } i=0 \\ 0 & \text{if } i>0\end{cases}$ where $\textrm{coind}_{*}^{G}(A) = \Hom_{\mathbb{Z}}(\mathbb{Z}[G], A)$ 
\end{corollary}

Take Shapiro lemma with $H = \left\{e\right\}$.

\begin{theorem}[Hilbert 90 (additive version)]
	Let $L/K$ be a finite Galois extension with Galois grp $G$.
	Then, we have that $H^i(G, L) = 
\begin{cases}
	K & \text{if } i = 0 \\
0 & \text{if } i>0 \\
\end{cases}
	$ 
\end{theorem}
\begin{proof}
By normal basis theorem, there exists $x \in L$ such that $\Orb_x$ under $G$ form a $K$-basis of $L$.
Now, consider $\Hom_{\mathbb{Z}}(\mathbb{Z}[G], K) \cong \mathbb{Z}[G] \otimes_{\mathbb{Z}} K \cong  K[G] \cong L$. 

Take $H = \left\{e\right\} \injto G$, $K$ as $H$-module.
We proved that $\textrm{coind}_{*}^G(K) = L$, so $H^i(G, L) = H^i(\left\{e\right\}, K)$ which is $K$ if $i=0$, trivial otherwise.  
\end{proof}

\begin{corollary}
Let $L/K$ be a cyclic extension with Galois group $G = \langle \sigma\rangle$.
Then, there exists an exact sequence 
\[
	0 \to K \to L \xrightarrow{\sigma-1} L \xrightarrow{\Tr} K \to 0
\]
\end{corollary}

\begin{proof}
By Hilbert 90, we have that $0 = H^i(G, L) = NL/(\sigma-1)L$.
We also have that $0 = H^2(G, L) = L^G/NL = L^G/\Tr(L) = K/\Tr(L)$ 
\end{proof}

\begin{corollary}
Let $L/K$ be a cyclic extension of degree $p$, where $p = \chara(K) >0$.
Then, $\exists x \in L$ such that $L = K(x)$ and $x^p - x \in K$.   
\end{corollary}
\begin{proof}
Take $1 \in L$ such that $\Tr(1) = p = 0$.
Thus, $\exists x \in L$ s.t. $\sigma(x) - x = 1$.
Thus, $\sigma(x) = x+ 1$, so $x$ is not fixed by $\sigma$ and $x  \not \in K$.
This, $L = K(x)$.
Then, expand $\sigma(x^p - x)$ and we get that $x^p - z \in K$. 
\end{proof}
 
Let $\alpha: R \to S$ be a (not necessarily commutative) ring homomorphism.
Suppose we have 
\begin{itemize}
	\item $S$-modules $P$ and $M$
	\item $R$-module $N$.
	\item An $R$-linear map $h: M \to N$   
\end{itemize}
\begin{theorem}
Under these conditions, there is a natural homomorphism $h_{*}: \Ext_S^i(P, M) \to \Ext_R^i(P, N)$, which is $R$-linear if $R$ is commutative.
\end{theorem}

\section{Class 7}
Suppose $G \circlearrowright A$ for an abelian group $A$, and there exists a homomorphism $\phi: G' \to G$.
Then, $A$ is a $G'$ module in the obvious manner.

Is there a connection between $H^i(G, A)$ and $H^i(G', A)$? 
Observe that there is a canonical inclusion $H^0(G, A) = A^G \injto A^{G'} = H^0(G', A)$.

We now prove the theorem proved in the last class.
Let $\alpha: R \to S$ be a ring homomorphism, and we have
\begin{enumerate}
	\item $S$-modules $P$ and $M$
	\item $R$-module $N$ 
	\item An $R$-linear map $h: M \to N$. 
\end{enumerate}

\begin{proof}
Let $M \to I^{\bullet}$ be an injective resolution of $M$ in $S-\catname{mod}$.
Let $I^{\bullet} \to J^{\bullet}$ be an injective resolution of $I^{\bullet}$.
Take the total complex,and look at maps 
\[
	\xymatrix{
	\Hom_S(P, M) \ar[r] & \Hom_S(P, I^{\bullet}) \ar[d]^{\alpha^*}\ar@{-->}[dl]\\
               \Hom_R(P, J^{\bullet}) & \Hom_R(P, I^{\bullet}) \ar[l] \\
       }
\]
Thus, there exists a map $\alpha^{*}: H^i(\Hom_S(P, I)) \to H^i(\Hom_R(P, J))$.
However, the LHS, resp. RHS are simply $\Ext_S^i(P, M) \to \Ext_R^i(P, M)$.  
\end{proof}

An application: Let $\phi: G' \to G$ be a group homomorphism and $A \in G-\catname{mod}$.
Let $B$ be a $G'$ module and let $h: A \to B$ be a $G'$ equivariant map.
Then, there exists a canonical group homomorphism $h_{*}$ 

\subsection{Class 8}

Last time our set up with: $\alpha: R \to S$ such that $S$ is a free $R$-module of finite rank $n$.
Let $M, P$ be $S$-modules.
Then, there exists a map 

\[
\xymatrix{
	\Ext_R^i(P, M) \ar@<0.5ex>[r]^{\text{Cores}^S_R} &  \Ext_S^i(P, M) \ar@<0.5ex>[l]^{\text{Res}_R^S}
}
\]

\begin{proposition}
	$\text{Cores}_R^S \circ \text{Res}_R^S : \Ext_S^i(P, M) \to \Ext_S^i(P, M)$ is multiplication by $n$ (rank of $S$ over $R$).
\end{proposition}

\begin{proof}
First, we consider the case for $i=0$
In this case, we have the diagram
\[
\xymatrix{
	\Hom_S(P, M) \ar@{^{(}->}[r]^{\text{Res}_R^S} & \Hom_R(P, M) \ar[r]^{\cong} & \Hom_S(P \otimes_R S, M) \ar[d]^{\cong} \\
				      \Hom_S(P, M) & \bigoplus_{i=1}^n \Hom_S(P, M) \ar[l]^{\text{Cores}_R^S}_{\triangle^{*}} & \Hom_S(P^n, M) \ar[l]^{\cong}
}
\]

such that the composition is the map $\sum_{I=1}^n f = nf$.

This gives a natural transformation $\eta: \Hom_S(P, \dashbk) \to \Hom_S(P, \dashbk)$ on $S-\catname{mod}$
such that $\eta(M)$ is multiplication by $n$ on $M$, so $\eta = n \cdot \id$
Thus, $n: \Ext_S^i(P, \dashbk) \to \Ext_S^i(P, \dashbk)$ is also equal to $n \cdot \id$ for all $i\geq 0$.   
\end{proof}

\begin{corollary}
	Let $G$ be a discrete group and $H \injto G$ a subgroup of finite index $n$.
	Let $A \in G-\catname{mod}$.
	Then, the composition $H^i(G, A) \xrightarrow{\text{Res}_H^G} H^i(H, A) \xrightarrow{\text{Cores}_H^G} H^i(G, A)$  
	is multiplication by $n = [G : H]$.
\end{corollary}
\begin{proof}
	Take $R = \mathbb{Z}[H]$ and $S = \mathbb{Z}[G]$, $P = \mathbb{Z}$, $M = A$.
\end{proof}
\begin{corollary}
	Suppose $G$ is a finite group of order $n$ and $A$ is a $G$-module.
	Then, $H^i(G, A)$ is annihilated by $n$ for all $i>0$. 
\end{corollary}
\begin{proof}
	Apply the previous corollary with $H = \left\{e\right\}$.
	Then, we get maps $H^i(G, A) \xrightarrow{\text{Res}} \to H^i(\left\{e\right\}, A) \xrightarrow{\text{Cores}} H^i(G, A)$  such that the composition is multiplication by $n$.
	Thus, $nH^i(G, A) = 0$ for all $i>0$.
\end{proof}
Let $G$ be a discrete group, $H$ a normal subgroup of finite index.
Let $A \in G-\catname{mod}$.
Then, we have maps
\[
\xymatrix{
	H^i(G, A) \ar[r]^{\text{Res}_H^G} \ar[dr]& H^i(H, A) \ar[d]^{\text{Cores}_H^G} \\
	H^i(G/H, A^H) \ar[u]^{\text{Inf}_H^G} \ar@{-->}[ur] & H^i(G, A) \\
}
\]

When $i =0$, we have the below diagram.
\[
\xymatrix{
	A^G \ar@{^{(}->}[r] & A^H \ar[d]^{\text{Cores}_H^G} \\
	A^G = \left(A^H\right)^{G/H} \ar@{^{(}->}[ur] \ar[u]^{\id} & A^G \\
}
\]
\end{document}


